\chapter{Dữ liệu quan hệ: Postgres, JPA, Flyway}
\label{ch:postgres}

Hai trong số các service, mỗi service sở hữu một cơ sở dữ liệu
PostgreSQL: auth-service sở hữu \code{ts\_auth}, course-service sở hữu
\code{ts\_course}. Chương này bàn ba tầng nằm giữa một phương thức
\code{@Service} và các byte trên đĩa --- JPA/Hibernate, connection pool,
và Flyway --- cùng quy tắc kiến trúc giải thích vì sao lại có tới
\emph{hai} cơ sở dữ liệu.

\section{Mỗi service một cơ sở dữ liệu}

\begin{decision}[hai server Postgres, không phải hai schema trong một]
Mỗi service sở hữu dữ liệu của mình một cách độc quyền; đường duy nhất
tới dữ liệu của service khác là qua API hoặc sự kiện của nó. Điều này
cấm kiểu thất bại kinh điển của cơ sở dữ liệu dùng chung --- service A
lặng lẽ join vào bảng của service B, trói schema hai bên với nhau mãi
mãi. Dự án đi xa hơn mức cần thiết: không chỉ cơ sở dữ liệu riêng mà
cả tiến trình server riêng, nên ngay cả kết nối chéo vô tình cũng bất
khả, và mỗi bên có thể được tinh chỉnh, sao lưu, hay làm hỏng độc lập.
Chi phí trên máy này: khoảng 256\,MiB cho mỗi server thêm. Trong một
tài khoản cloud, cùng mức cô lập đó thường là hai cơ sở dữ liệu trên
một instance RDS được quản lý --- cô lập bằng thông tin đăng nhập thay
vì bằng tiến trình --- vì instance được quản lý tính tiền theo server.
\end{decision}

\section{JPA và Hibernate}

JPA là đặc tả (các annotation như \code{@Entity},
\code{@OneToMany}, \code{EntityManager}); Hibernate là bản cài đặt mà
Spring Boot tự cấu hình. Spring Data JPA nằm bên trên, suy ra truy vấn
từ tên phương thức repository (\code{findByCourseIdAndStudentId}) và
sinh phần mã lặp.

Cấu hình từ \code{configurations/course-service.yml}:

\begin{lstlisting}[language=yaml]
spring:
  datasource:
    url: jdbc:postgresql://localhost:5432/ts_course   (*@\co{1}@*)
    username: postgres
    password: root                                    (*@\co{2}@*)
  jpa:
    hibernate:
      ddl-auto: none                                  (*@\co{3}@*)
    show-sql: true                                    (*@\co{4}@*)
  flyway:
    enabled: true                                     (*@\co{5}@*)
\end{lstlisting}

\begin{description}
  \item[\co{1}] Mặc định khi dev. Compose ghi đè host thành
  \code{postgres-course}; Kubernetes ghi đè y hệt qua
  \code{SPRING\_DATASOURCE\_URL}. Để ý Service k8s được cố ý \emph{đặt
  tên} \code{postgres-course} để chuỗi ghi đè giống nhau ở cả hai môi
  trường.
  \item[\co{2}] Thông tin đăng nhập dạng văn bản thuần chỉ chấp nhận
  được vì file này là \emph{mặc định khi dev}; trong cụm,
  \code{SPRING\_DATASOURCE\_USERNAME}/\code{PASSWORD} tới từ một Secret
  Kubernetes (Chương~16) và có ưu tiên cao hơn.
  \item[\co{3}] Dòng quan trọng nhất trong khối ---
  Mục~\ref{sec:ddlauto}.
  \item[\co{4}] Mọi câu SQL đều được log --- một công cụ dạy học và một
  máy dò truy vấn N+1 khi dev; tắt đi ở nơi dung lượng log tốn tiền.
  \item[\co{5}] Quản lý schema thuộc về Flyway, không phải Hibernate.
\end{description}

\section{Vì sao \code{ddl-auto: none}}
\label{sec:ddlauto}

Hibernate có thể sinh schema từ entity (\code{ddl-auto: update}), và
tutorial nào cũng dùng nó. Hệ thống production thì không, vì
\code{update} là cái bánh cóc một chiều không có lịch sử: nó thêm cột
nhưng không bao giờ xóa hay đổi tên một cách an toàn; cơ sở dữ liệu của
hai lập trình viên lệch nhau âm thầm; không có bản ghi nào về cái gì đổi
khi nào, và không có cách nào xem xét một thay đổi schema trước khi nó
chạm tới production.

\section{Flyway: migration dưới dạng code}

Flyway coi schema là một chuỗi có thứ tự các script SQL bất biến, mỗi
script áp dụng đúng một lần, được theo dõi trong một bảng
(\code{flyway\_schema\_history}) ngay trong cơ sở dữ liệu đích:

\begin{lstlisting}
auth-service/src/main/resources/db/migration/
  V1__create_users.sql
  V2__create_refresh_tokens.sql
  V3__add_admin_user_invitations.sql
  V4__seed_default_admin.sql
\end{lstlisting}

Lúc khởi động, Flyway so bảng lịch sử với các script trên classpath và
áp dụng những script còn thiếu, theo thứ tự phiên bản, trước khi
Hibernate kịp khởi tạo. Các quy tắc làm nó hoạt động: không bao giờ sửa
một migration đã áp dụng (mỗi dòng lưu một checksum --- sửa sẽ bị phát
hiện và từ chối); chỉ thêm phiên bản mới; và viết migration tương thích
với phiên bản ứng dụng \emph{trước đó} khi bạn triển khai cuốn chiếu
(Chương~15): thêm-rồi-chuyển-rồi-xóa, không bao giờ đổi tên trong một
bước.

\begin{pitfall}[chạy tốt ở máy mình, CrashLoopBackOff trong cụm]
Triệu chứng: service khởi động được trên máy bạn nhưng pod của nó khởi
động lại không ngừng; log cho thấy Flyway lệch checksum hoặc một
migration thất bại. Nguyên nhân: cơ sở dữ liệu cục bộ đã được migrate
bằng một script cũ hơn (hoặc đã bị sửa) --- bảng lịch sử trên laptop của
bạn và của cụm không khớp nhau. Đây cũng là lý do startup probe của
Kubernetes (Chương~17) nhắm vào một endpoint chỉ trả lời sau khi Spring
context đã lên: ``migration đã áp dụng'' là một phần của nghĩa
\emph{sẵn sàng}. Sửa bằng cách thêm một migration hiệu chỉnh, không bao
giờ bằng cách sửa lịch sử.
\end{pitfall}

\section{Connection pool}

Mở một kết nối Postgres tốn một bắt tay TCP, xác thực, và một lần fork
tiến trình server --- quá đắt cho mỗi request. HikariCP (pool mặc định
của Spring Boot) giữ một nhóm nhỏ kết nối đang mở để request mượn rồi
trả. Tối đa mặc định: 10.

\begin{hardware}
Truyền miệng về kích thước pool nói ``nhiều kết nối hơn, thông lượng cao
hơn''; thực tế trên server 2 nhân thì ngược lại --- Postgres làm việc
\emph{theo từng kết nối}, và công thức thực dụng
(\(\approx\) số nhân $\times$ 2) cho thấy mặc định 10 đã là hào phóng
với hai service chia nhau hai nhân. Nếu Postgres có lúc báo
\code{too many clients}, hãy hạ pool xuống trước khi nâng
\code{max\_connections}.
\end{hardware}

\section{Ngoài dự án}

Phiên bản cloud được quản lý tương đương với thiết lập của chương này
--- AWS RDS Postgres với cùng luồng Flyway --- chỉ thay đổi URL và mô
hình tin cậy:

\begin{lstlisting}[language=yaml]
spring:
  datasource:
    url: jdbc:postgresql://app-db.xxxxx.eu-central-1.rds.amazonaws.com:5432/ts_course
    username: ${DB_USER}        # from AWS Secrets Manager
    password: ${DB_PASSWORD}
\end{lstlisting}

Sao lưu, sao chép (replication) và failover thành việc của nhà cung
cấp --- đó là lý do cảnh báo của Chương~19 về trạng thái tự host lại
quan trọng. Biến thể công nghiệp thứ hai đáng biết: chạy Flyway như một
bước CI riêng (\code{flyway migrate} vào đích) thay vì lúc ứng dụng khởi
động, để thay đổi schema triển khai độc lập với code và một migration
chậm không ăn hết ngân sách của startup probe.

\begin{quiz}
\begin{enumerate}
  \item Nêu hai kiểu thất bại cụ thể của \code{ddl-auto: update} trong
  môi trường làm việc nhóm mà \code{none} + Flyway ngăn được.
  \item Vì sao không bao giờ được sửa một migration đã áp dụng, và
  Flyway dùng gì để phát hiện?
  \item Bạn cần đổi tên một cột mà không có downtime dưới triển khai
  cuốn chiếu. Phác thảo chuỗi nhiều bản phát hành.
  \item Bảo vệ và phản bác quyết định hai server: mỗi phía một lập
  luận, và bối cảnh nào thì phía nào thắng.
\end{enumerate}
\end{quiz}
